\documentclass[10pt,a4paper]{amsart}
\usepackage{amsmath,amssymb,amsthm,graphicx}
% Retained minimal preamble for the Stacks Project chapter varieties.tex.
% Verbatim excerpt of the chapter preamble: Xy-pic setup (native preamble lines 14-18), the native theorem
% declarations with the mathematical macros (native preamble lines 154-217), and the project own no-op
% definitions of the reference/slogan/history display environments (native wrapper.tex lines 1-3) that the
% counted statement uses for its slogan. Copied mechanically; nothing added, removed or redefined.
% For commutative diagrams we use Xy-pic
\usepackage[all]{xy}

% We use 2cell for 2-commutative diagrams.
\xyoption{2cell}
% Theorem environments.
%
\theoremstyle{plain}
\newtheorem{theorem}[subsection]{Theorem}
\newtheorem{proposition}[subsection]{Proposition}
\newtheorem{lemma}[subsection]{Lemma}

\theoremstyle{definition}
\newtheorem{definition}[subsection]{Definition}
\newtheorem{example}[subsection]{Example}
\newtheorem{exercise}[subsection]{Exercise}
\newtheorem{situation}[subsection]{Situation}

\theoremstyle{remark}
\newtheorem{remark}[subsection]{Remark}
\newtheorem{remarks}[subsection]{Remarks}

\numberwithin{equation}{subsection}

% Macros
%
\def\lim{\mathop{\mathrm{lim}}\nolimits}
\def\colim{\mathop{\mathrm{colim}}\nolimits}
\def\Spec{\mathop{\mathrm{Spec}}}
\def\Hom{\mathop{\mathrm{Hom}}\nolimits}
\def\Ext{\mathop{\mathrm{Ext}}\nolimits}
\def\SheafHom{\mathop{\mathcal{H}\!\mathit{om}}\nolimits}
\def\SheafExt{\mathop{\mathcal{E}\!\mathit{xt}}\nolimits}
\def\Sch{\mathit{Sch}}
\def\Mor{\mathop{\mathrm{Mor}}\nolimits}
\def\Ob{\mathop{\mathrm{Ob}}\nolimits}
\def\Sh{\mathop{\mathit{Sh}}\nolimits}
\def\NL{\mathop{N\!L}\nolimits}
\def\CH{\mathop{\mathrm{CH}}\nolimits}
\def\proetale{{pro\text{-}\acute{e}tale}}
\def\etale{{\acute{e}tale}}
\def\QCoh{\mathit{QCoh}}
\def\Ker{\mathop{\mathrm{Ker}}}
\def\Im{\mathop{\mathrm{Im}}}
\def\Coker{\mathop{\mathrm{Coker}}}
\def\Coim{\mathop{\mathrm{Coim}}}

% Boxtimes
%
\DeclareMathSymbol{\boxtimes}{\mathbin}{AMSa}{"02}

%
% Macros for moduli stacks/spaces
%
\def\QCohstack{\mathcal{QC}\!\mathit{oh}}
\def\Cohstack{\mathcal{C}\!\mathit{oh}}
\def\Spacesstack{\mathcal{S}\!\mathit{paces}}
\def\Quotfunctor{\mathrm{Quot}}
\def\Hilbfunctor{\mathrm{Hilb}}
\def\Curvesstack{\mathcal{C}\!\mathit{urves}}
\def\Polarizedstack{\mathcal{P}\!\mathit{olarized}}
\def\Complexesstack{\mathcal{C}\!\mathit{omplexes}}
% \Pic is the operator that assigns to X its picard group, usage \Pic(X)
% \Picardstack_{X/B} denotes the Picard stack of X over B
% \Picardfunctor_{X/B} denotes the Picard functor of X over B
\def\Pic{\mathop{\mathrm{Pic}}\nolimits}
\def\Picardstack{\mathcal{P}\!\mathit{ic}}
\def\Picardfunctor{\mathrm{Pic}}
\def\Deformationcategory{\mathcal{D}\!\mathit{ef}}
\newenvironment{reference}{}{}
\newenvironment{slogan}{}{}
\newenvironment{history}{}{}
\usepackage{cleveref}
\crefname{originallemma}{original lemma}{original lemmas}
\Crefname{originallemma}{Original lemma}{Original lemmas}
\makeatletter
\newlabel{algebra-lemma-characterize-smooth-over-field}{{00TT}{source}{}{sourcepointer-1}{}}
\newlabel{algebra-lemma-characterize-smooth-over-field@cref}{{[originallemma][0][]00TT}{[1][0][]source}{}{}{}}
\makeatother
\crefname{originallemma}{original lemma}{original lemmas}
\Crefname{originallemma}{Original lemma}{Original lemmas}
\makeatletter
\newlabel{algebra-lemma-make-separable}{{030R}{source}{}{sourcepointer-2}{}}
\newlabel{algebra-lemma-make-separable@cref}{{[originallemma][0][]030R}{[1][0][]source}{}{}{}}
\makeatother
\crefname{originallemma}{original lemma}{original lemmas}
\Crefname{originallemma}{Original lemma}{Original lemmas}
\makeatletter
\newlabel{algebra-lemma-p-ring-map}{{0BRA}{source}{}{sourcepointer-3}{}}
\newlabel{algebra-lemma-p-ring-map@cref}{{[originallemma][0][]0BRA}{[1][0][]source}{}{}{}}
\makeatother
\crefname{originallemma}{original lemma}{original lemmas}
\Crefname{originallemma}{Original lemma}{Original lemmas}
\makeatletter
\newlabel{algebra-lemma-separable-smooth}{{00TV}{source}{}{sourcepointer-4}{}}
\newlabel{algebra-lemma-separable-smooth@cref}{{[originallemma][0][]00TV}{[1][0][]source}{}{}{}}
\makeatother
\crefname{originallemma}{original lemma}{original lemmas}
\Crefname{originallemma}{Original lemma}{Original lemmas}
\makeatletter
\newlabel{algebra-lemma-smooth-field-change-local}{{02UQ}{source}{}{sourcepointer-5}{}}
\newlabel{algebra-lemma-smooth-field-change-local@cref}{{[originallemma][0][]02UQ}{[1][0][]source}{}{}{}}
\makeatother
\crefname{originaldefinition}{original definition}{original definitions}
\Crefname{originaldefinition}{Original definition}{Original definitions}
\makeatletter
\newlabel{morphisms-definition-smooth}{{01V5}{source}{}{sourcepointer-6}{}}
\newlabel{morphisms-definition-smooth@cref}{{[originaldefinition][0][]01V5}{[1][0][]source}{}{}{}}
\makeatother
\title{Geometric regularity at a point of a locally finite type scheme over a field is equivalent to smoothness there}
\author{The Stacks Project Authors}
\date{stacks-project:varieties.tex:local-snapshot-20261011; GNU FDL 1.2 or later}
\begin{document}
\maketitle
\noindent\textit{Editorial note.} All delivered mathematical text is the author's own, copied from the pinned Stacks Project chapter file varieties.tex; no statement or proof step was rewritten, reordered or removed.\par\medskip
\noindent\textit{Editorial note.} The proof is complete as written and ends with the descent of smoothness along the purely inseparable extension; no step is referred to the reader.\par\medskip
\noindent\textit{Editorial note.} The counted spans carry 6 navigation stubs and no bibliography citation.\par\medskip
\noindent\textit{Editorial note.} The preamble is the chapter's retained minimal preamble (sources/stacks-varieties/retained-preamble.tex lines 1-77), a verbatim excerpt of the native preamble.\par\medskip
\noindent\textit{Editorial note.} The label lemma-geometrically-regular-smooth is defined only in varieties.tex, so no cross-chapter restatement had to be weighed for this card.\par\medskip
\noindent\textit{Allowed standard prerequisites (editorial).}\begin{itemize}
\item Smooth algebras over a field are regular and smoothness is stable under field extension. Algebra, Lemma algebra-lemma-characterize-smooth-over-field, invoked at proof lines 2560-2563. Supplies the easy direction of the equivalence.
\item Making a finitely generated field extension separable by a finite purely inseparable base change. Algebra, Lemma algebra-lemma-make-separable, invoked at proof lines 2569-2572. Supplies the finite purely inseparable extension k' over which the residue field becomes separable.
\item Uniqueness of the prime above for a purely inseparable extension. Algebra, Lemma algebra-lemma-p-ring-map, invoked at proof lines 2573-2575. Supplies that the residue field at the lifted prime is the compositum k'K.
\item Separable residue field plus regularity implies smoothness, and smoothness descends along a purely inseparable field extension. Algebra, Lemmas algebra-lemma-separable-smooth and algebra-lemma-smooth-field-change-local, invoked at proof lines 2578-2582. Supplies the two algebraic steps connecting regularity to smoothness over k.
\end{itemize}\medskip
\ifdefined\hypertarget\hypertarget{sourcepointer-1}{}\fi
\noindent\textit{Source reference.} Original lemma 00TT: sources/stacks-chow/tags.txt, source lines 1087--1087. Pointer outside the counted span, invoked once for context; the load-bearing argument of the counted statement is the author's proof reproduced in full in the proof span.\par\medskip
\ifdefined\hypertarget\hypertarget{sourcepointer-2}{}\fi
\noindent\textit{Source reference.} Original lemma 030R: sources/stacks-chow/tags.txt, source lines 3958--3958. Pointer outside the counted span, invoked once for context; the load-bearing argument of the counted statement is the author's proof reproduced in full in the proof span.\par\medskip
\ifdefined\hypertarget\hypertarget{sourcepointer-3}{}\fi
\noindent\textit{Source reference.} Original lemma 0BRA: sources/stacks-chow/tags.txt, source lines 14763--14763. Pointer outside the counted span, invoked once for context; the load-bearing argument of the counted statement is the author's proof reproduced in full in the proof span.\par\medskip
\ifdefined\hypertarget\hypertarget{sourcepointer-4}{}\fi
\noindent\textit{Source reference.} Original lemma 00TV: sources/stacks-chow/tags.txt, source lines 1089--1089. Pointer outside the counted span, invoked once for context; the load-bearing argument of the counted statement is the author's proof reproduced in full in the proof span.\par\medskip
\ifdefined\hypertarget\hypertarget{sourcepointer-5}{}\fi
\noindent\textit{Source reference.} Original lemma 02UQ: sources/stacks-chow/tags.txt, source lines 3741--3741. Pointer outside the counted span, invoked once for context; the load-bearing argument of the counted statement is the author's proof reproduced in full in the proof span.\par\medskip
\ifdefined\hypertarget\hypertarget{sourcepointer-6}{}\fi
\noindent\textit{Source reference.} Original definition 01V5: sources/stacks-chow/tags.txt, source lines 2434--2434. Pointer outside the counted span, invoked once for context; the load-bearing argument of the counted statement is the author's proof reproduced in full in the proof span.\par\medskip
\begin{lemma}
\label{lemma-geometrically-regular-smooth}
Let $k$ be a field.
Let $X$ be a scheme locally of finite type over $k$.
Let $x \in X$.
Then $X$ is geometrically regular at $x$ if and only if $X \to \Spec(k)$
is smooth at $x$ (Morphisms, Definition \ref{morphisms-definition-smooth}).
\end{lemma}

\begin{proof}
The question is local around $x$,
hence we may assume that $X = \Spec(A)$
for some finite type $k$-algebra.
Let $x$ correspond to the prime $\mathfrak p$.

\medskip\noindent
If $A$ is smooth over $k$ at $\mathfrak p$, then we may localize $A$
and assume that $A$ is smooth over $k$. In this case $k' \otimes_k A$
is smooth over $k'$ for all extension fields $k'/k$, and each of
these Noetherian rings is regular by
Algebra, Lemma \ref{algebra-lemma-characterize-smooth-over-field}.

\medskip\noindent
Assume $X$ is geometrically regular at $x$.
Consider the residue field $K := \kappa(x) = \kappa(\mathfrak p)$ of $x$.
It is a finitely generated extension of $k$.
By Algebra, Lemma \ref{algebra-lemma-make-separable}
there exists a finite purely inseparable
extension $k'/k$ such that the compositum
$k'K$ is a separable field extension of $k'$.
Let $\mathfrak p' \subset A' = k' \otimes_k A$ be a prime ideal
lying over $\mathfrak p$. It is the unique prime lying over $\mathfrak p$, see
Algebra, Lemma \ref{algebra-lemma-p-ring-map}.
Hence the residue field $K' := \kappa(\mathfrak p')$
is the compositum $k'K$. By assumption the local ring
$(A')_{\mathfrak p'}$ is regular. Hence by
Algebra, Lemma \ref{algebra-lemma-separable-smooth}
we see that $k' \to A'$ is smooth at $\mathfrak p'$.
This in turn implies that $k \to A$ is smooth at $\mathfrak p$ by
Algebra, Lemma \ref{algebra-lemma-smooth-field-change-local}.
The lemma is proved.
\end{proof}

\end{document}
